\documentclass[11pt,reqno]{amsart}
\usepackage[T1]{fontenc}\usepackage[utf8]{inputenc}\usepackage{lmodern}
\usepackage[a4paper,margin=23mm]{geometry}\usepackage{amsmath,amssymb,mathtools,microtype,xurl}
\usepackage[hidelinks]{hyperref}
\newtheorem{theorem}{Theorem}\newtheorem{proposition}[theorem]{Proposition}
\title[Exact two-channel and reciprocal-source comparison]{An Exact Two-Channel Transform and a Reciprocal-Source Obstruction\\for the Erd\H{o}s--Straus Equation}
\author{The Clankers}\date{19 September 2026}
\begin{document}
\begin{abstract}
On every negative-character source in the square-source ES system, exterior and
middle hits require opposite original divisor characters. A common p-coloured
pair source has exact fibres of size $2\tau_{\mathrm{div}}(h)$, whose M\"obius transform recovers
the full joint count. The character forcing new source factors is annihilated in
that count. We give a finite local obstruction to immediate reciprocal-source
repair, and an original non-square exchange that crosses the channels elsewhere
at the same prime. No universal ES nonvanishing result is asserted.
\end{abstract}\maketitle

Let $p\equiv1\pmod4$ be prime, $p/4<a<p/2$, $R=4a-p$, and $u\mid a^2$.
The original labelled gates are
\[
 E:R\mid4u+1,\qquad M:R\mid4u+p.
\]
With $d=\gcd(a,u)$ and $(h,r,s)=(d^2/u,u/d,a/d)$ their ordered returns are
\[
 E:(a,hs\kappa,phr\kappa),\quad\kappa=(pr+s)/R;
\qquad M:(a,phs\lambda,phr\lambda),\quad\lambda=(r+s)/R.  \tag{1}
\]
Their ordered inverses are $a^2/(Ry-pa)$ and $pa^2/(Ry-pa)$ respectively.
These are inherited Type I/II coordinates [1]. Every coordinate is retained.

\begin{proposition}[Channel characters]
For every original divisor, $(u/R)=(u/p)$. An E hit requires $(u/p)=-1$;
an M hit requires $(u/p)=-(a/p)$. Thus on a negative source $(a/p)=-1$,
unchanged-$u$ and rational-square-ratio transports cannot cross E and M.
\end{proposition}
\begin{proof}
For odd $\ell\mid a$, reciprocity and $R\equiv-p\pmod\ell$ give
$(\ell/R)=(-1/\ell)(-p/\ell)=(\ell/p)$. The case 2 follows from
$R\equiv-p\pmod8$ when $2\mid a$. Multiply over the original exponents.
The two gate residues are $-1/4$ and $-p/4$, and $(p/R)=(a/p)$.
\end{proof}
In particular, for $a=(p+\sigma_q qt^2)/4$ with $q$ an external nonresidue
prime and a permitted hard-prime square source, the two required signs are
opposite. This is a restriction on successful returns, not a sufficiency test.

For $\gcd(pN,R)=1$, $p\nmid N$, retain the single ordered pair source
\[
 \mathcal P_R(N)=\{(b,c):b,c\mid pN,\ R\mid b+c\},\quad T_R(N)=|\mathcal P_R(N)|.
\]
Let $F_R(N)$ count ordered coprime pairs $rs\mid N$ satisfying $R\mid r+s$
(label M) or $R\mid pr+s$ (label E), as a disjoint tagged union.

\begin{theorem}[Exact fibres and joint count]
At a primitive pair put $h=N/(rs)$. Its complete fibre is
\[
 M:\ (p^\varepsilon dr,p^\varepsilon ds),\quad d\mid h,\ \varepsilon=0,1;
\]
\[
 E:\ (pdr,ds)\ \hbox{or}\ (ds,pdr),\quad d\mid h.
\]
It has size $2\tau_{\mathrm{div}}(h)$. Consequently
\[
 \boxed{F_R(N)=\frac12\sum_{d\mid N}\mu_{\mathrm M}(d)T_R(N/d).}      \tag{2}
\]
At the actual shell $N=a$, $F_R(a)=|E_a|+|M_a|$, so the complete first-half count is
\[
 \boxed{\mathcal C_p=\frac12\sum_{p/4<a<p/2}\sum_{d\mid a}
 \mu_{\mathrm M}(d)T_{4a-p}(a/d).}                                   \tag{3}
\]
\end{theorem}
\begin{proof}
Divide $(b,c)$ by its gcd. If neither coprime part contains p, remove the common
p-bit and recover the first displayed fibre. If exactly one contains p, remove
it, retain its side, and recover the second. The lcm condition gives $d\mid h$.
These operations are inverse. Thus $T_R(N)/2=\sum_{d\mid N}F_R(N/d)$;
M\"obius inversion follows from $\sum_{d\mid n}\mu_{\mathrm M}(d)=\mathbf1_{n=1}$.
At $N=a$, the inverse $u=hr^2$ gives exactly the original states. Auxiliary
$N=a/d$ are divisor-pair sources, not shells satisfying $4N=p+R$.
\end{proof}

For existence, the exact fibres also give
$2F_R(a)\le T_R(a)\le2\tau_{\mathrm{div}}(a)F_R(a)$. Hence (3) is positive exactly when
$\sum_{p/4<a<p/2}T_{4a-p}(a)>0$. One can force an original pair first and
normalize afterward; termwise positivity of the M\"obius summands is unnecessary.

The two p-colours give the matrix
\[
 \begin{pmatrix}B_0&B_1\\B_1&B_0\end{pmatrix},\qquad
 B_0=\#\{b,c\mid N:R\mid b+c\},\quad B_1=\#\{b,c\mid N:R\mid b+pc\}.
\]
It need not be positive semidefinite. At $p=193,R=7,a=50$, its entries are
$B_0=4,B_1=6$; twenty ordered pairs normalize to four original states, not ten.

Character orthogonality and the local Fej\'er difference give
\[
 F_R(N)=\frac1{2\varphi(R)}\sum_{\chi\bmod R}
 \chi(-1)|1+\chi(p)|^2
 \prod_{q^e\parallel N}\sum_{j=-e}^e\chi(q)^j.             \tag{4}
\]
Indeed $T_R(N)$ has the same expression with the last product replaced by
$|\sum_{d\mid N}\chi(d)|^2$; subtracting the previous prime exponent gives
$|1+z+\cdots+z^e|^2-|1+z+\cdots+z^{e-1}|^2=\sum_{j=-e}^ez^j$.
At every negative square source the transported Jacobi character satisfies
$\chi_R(p)=-1$, so its entire term is killed by $|1+\chi_R(p)|^2$.
Thus the character forcing factors does not itself force the combined target.
The other signed terms remain an arithmetic obligation.

\begin{proposition}[Reciprocal source and exact inverse]
For $D\equiv3\pmod4$, let $B=(pD+1)/4$. A divisor $w\mid B^2$ with
$D\mid w+B$ normalizes as $B=hrs,w=hr^2$, $\gcd(r,s)=1$ and
$r+s=D\lambda$. After recording the swap making $r<s$, it returns
\[
 (x,y,z)=(hr\lambda,hs\lambda,pB),\qquad p/4<x<p/2.
\]
It is the E state $(h,r,s_{\rm original},\kappa)=(h,r,\lambda,s)$.
The inverse is $D=(4u+1)/R$, $B=hr\kappa$, with the two source divisors
$w=hr^2$ and $h\kappa^2$.
\end{proposition}
\begin{proof}
Use $4hrs=pD+1$ and $r+s=D\lambda$. The original residual satisfies
$(4x-p)s=pr+\lambda>0$. Moreover
$(4hrs-1)(s-r)-(r+s)=2s(2hr(s-r)-1)>0$, proving the strict upper bound.
The displayed identities solve back for D, B and both orientations. The symmetric
master equation is inherited [2], not a new existence assertion.
\end{proof}
Every first-half E state has $D\le\lfloor(p^2+6p+13)/12\rfloor$, by
$u\le a^2$ and maximizing $(4a^2+1)/(4a-p)$ at residual 3.

\begin{proposition}[Overlap supplies only the old endpoint]
For $p\equiv1\pmod8$, $0<D<3p$, put $A=(p+D)/4$, $B=(pD+1)/4$.
A nonresidue prime dividing both must be $\ell\equiv3\pmod4$ with
$\ell\mid p+1$. Conversely any such factor gives a common nonresidue overlap
for some $D<p$.
\end{proposition}
\begin{proof}
$\gcd(A,B)=\gcd(A,(D^2-1)/4)$. At a common prime $D=\pm1$ and $p=-D$;
reciprocity rules out the sign $p=1$ and requires $\ell\equiv3\pmod4$ otherwise.
For the converse take the least such $\ell$. Since $(p+1)/2\equiv1\pmod4$,
$\ell^2\le(p+1)/2$; then $D=2\ell+1<p$ gives the overlap. The original
endpoint E return is $R=\ell$, $a=u=(p+\ell)/4$, $\kappa=(p+1)/\ell$.
\end{proof}

At the certified prime $p=12889$, consider $q=43,61$ and all their valid square
ports $D=\sigma_q qt^2<3p$: respectively $t=1,3,\ldots,29$ and
$t=1,3,\ldots,13$. The complete finite certificate exhausts the 246 original
divisors in the 22 direct sources (the inherited Turn 4 fixture [3]) and the
600 original divisors in all 22 reciprocal sources. All direct E/M and all
reciprocal sum gates are empty. This is 44 boxes and 66 gate families, not a
whole-graph or ES counterexample. At the canonical ports the reciprocal sources
are $127\cdot1091$ and $2^3\cdot73709$, with no nonresidue factors. The two-cycle
therefore remains closed under that canonical-only reciprocal augmentation.
A reciprocal nonresidue factor is not sufficient either: $B_{387}=37\cdot33703$
has two, but its gate is empty.

At the same p, the genuine crossing occurs at $R=31$, $a=3230$:
\[
 u_M=25,\quad u_E=85,\quad \frac{u_E}{u_M}=\frac{17}{5}\equiv p^{-1}\pmod{31}.
\]
This ratio has character $-1$ and stays within the actual prime-exponent box.
The returns are $(3230,174852174,1353345)$ and $(3230,1346910,456850605)$.
The reciprocal source at $q=61,t=11$ has factorization
$2\cdot19\cdot31\cdot37\cdot1637$ and supplies the actual nonresidue factor 31
leading to that occupied source. This is a path, not a universal path theorem.

The next target is a positive lower bound for the complete original sum (3).
No such bound follows from (4)'s annihilated source character or from the formal
reciprocal inverse. The standard-library certificate and separate implementation
state their finite bounds and retain the original order, characters and fibres.
No ES resolution, leastness, independent review, Lean build or historical
priority is claimed.

\begin{thebibliography}{9}
\bibitem{1} C. Elsholtz and T. Tao (2013), \emph{Counting the number of solutions to the
Erd\H{o}s--Straus equation on unit fractions}, arXiv:1107.1010, \S2.
\bibitem{2} C. Bryan (2026), \emph{FAB-DUAL-DESCENT-SYSTEM.md}, CENTL revision
\texttt{42508684c4386b575d7b52e763565345a76bcb9e}, \S\S1--4.
\bibitem{3} The Clankers (2026), \emph{Turn 4: mixed square-source finite reduction},
PR19 revision \texttt{4b6bb0f187fa44bb62fc8c51e996f05dfc90e2d7}, \S8.
\end{thebibliography}
\end{document}
